% Historical (v1) numbers below come from hash-pinned summary reports, not from original run
% receipts; see Appendix F for the provenance gap. v2 numbers come from receipts under paper/evidence/.
\section{Three ways the evaluation misled us}
\label{sec:lies}

Each subsection states what we believed, what the record shows, and the correct statement. The
three cases differ in kind. The first is a descriptive comparison between systems that differ in
more than one way. The second is a provenance failure. Only the third has a paired intervention
that isolates a cause. Numbers for the first two come from the earlier benchmark generation (v1:
18 splits, 5{,}800 frames, same generator lineage, a1-light convention throughout) and from public
real-photo sets. They are hash-pinned historical summaries, not original run receipts; Appendix~\ref{app:receipts}
records the gap. Numbers for the third come from the paired parity probes and from benchmark v2.
The probe, baseline and holdout-check result summaries are in the evidence index (Appendix~\ref{app:receipts}); per-image prediction files are not.

\subsection{A leaderboard-selected pipeline scores below its own member}
\label{sec:router}

The system we had shipped routed each image, by the detected board's apparent angle, to one of
three members: a full-frame recogniser initialised from ChessQueries (the \emph{raw member}) and
two recognisers that operate on a rectified board crop (the \emph{warp-high} and \emph{warp-low}
members). The router and its thresholds had been chosen on a private leaderboard. On our 18
controlled splits the shipped ensemble scores below its own raw member on 16, by up to 32.3
exact-board points on one piece family. The raw member alone reaches 100.0 on twelve splits
(Figure~\ref{fig:router}). Section~\ref{sec:cost} reports the latency of both systems.

Three further comparisons point the same way. Each carries a caveat that makes it a development
number, not a replication. SLCC test: 52.45 against 74.73, on 368 boards from 13 videos, and the
raw member had SLCC training data in its mix. CVChess: 1.42 against 98.86, on 352 photographs of
only 88 distinct positions. A held-out block of the training scenes: 99.30 against 100.00, on
held-out move blocks of the same games and scenes as training. On SLCC, above the routing
threshold the shipped ensemble drops the raw member and uses only the two warp members. Run alone
on all 368 boards, those members produce eight or more wrong squares on 234 and 236 boards. That
is consistent with failures in the detector and rectification stages on broadcast crops. We ran no
oracle-corner ablation on that cohort, so localisation and recognition failures are not separated.

Routing, member composition and training lineage all differ between the shipped ensemble and the
raw member, so this is not an isolated routing effect. We also have no intervention on leaderboard
exposure, so the reversal cannot be attributed to adaptive reuse rather than to distribution shift
between the leaderboard and these cohorts. What the record supports is this: a pipeline selected
% pooled v1: n-weighted over paper/evidence/reports/historical_pipeline.csv, 5,800 frames
on one private cohort scored below its own simpler member on the pooled v1 benchmark (91.8 against
98.5 exact over 5{,}800 frames) and on each of the three further cohorts. The two exceptions among
the 18 splits are the two lowest elevation bands, where the pipeline scores 88.8 against 88.6 and
98.0 against 97.0.

\subsection{Two holdouts contaminated by construction}
\label{sec:contamination}

An identifier-level leak check over the training mixes and every named holdout found that two
holdouts derived from ChessRender360 (1{,}635 and 1{,}263 images) shared image identifiers with
the mixes trained on them: 1{,}446 and 1{,}128 identifiers, 2{,}398 distinct, 2{,}574 counting both
sets (5{,}148 mix-by-identifier events over two mixes). Both holdouts were subsets of the same ten
thousand rectified images. The check reports
\texttt{clean: false}. The documented model-selection replay does not use those two holdouts, so
the recorded selection decisions were not influenced by them. That is not a proof that no earlier
or unrecorded decision was. A set can be named a holdout and still be training data. A check that
runs on identifiers and positions, not on directory names, is what caught it.

\subsection{A convention confound mistaken for a domain shift}
\label{sec:parity}

The ChessQueries release checkpoint produces zero exact boards on every one of the 18 v1 splits.
Its per-square accuracy of 52.7--56.7 is 4.6--9.0 points \emph{below} a constant all-empty
prediction, which also produces zero exact boards on these positions. The model places pieces on
the wrong squares; it does not default to empty. The same checkpoint through the same code path
scores 99.62 on ChessReD and 98.54 on ChessCog, which is also rendered. The failure appeared specific to
our rendering lineage. A crop diagnostic fed the checkpoint the first 150 boards of each of three
splits (450 in all) four ways: full frame, bounding-box crop, square-padded crop, and a perspective
warp to a square. None recovered an exact board under our convention, and per-square accuracy
moved by at most 1.2 points from full frame. At that point our own notes read ``the shift is the
piece and board appearance itself.''

It was not. Our renderer, following the convention of the data it had been matched to, makes a1 a
\emph{light} square. The laws of chess place a dark square at a1 \cite{fide}. Every public
chess-recognition source we could resolve follows that convention (Table~\ref{tab:conventions}).
The audit checks one representative image per source; the SLCC sample has no printed coordinates,
so its corner is recorded as unresolved. A preregistered paired probe changed exactly that one
renderer expression and nothing else, with byte-identical controls. On 192 positions in two
disjoint waves the checkpoint went from 0/192 to 18/192 exact ($+9.38$ points, 95\% paired bootstrap
$[5.73,13.54]$) and from 54.36 to 85.24 per-square ($+30.88$, $[28.91,32.76]$), removing 67.65\% of
wrong squares. Replicated on 288 positions across six other piece families,
the effect is 0/288 to 50/288 exact and 55.31 to 87.49 per-square. Swapping the model's learned
two-class query-colour indices \cite{chessqueries} in memory recovers 20.7 per-square points on a1-light images
but no exact board, and costs 2.9 points on a1-dark images. This is consistent with the colour
prior contributing to the failure; it does not isolate the mechanism. On the
development partition of benchmark v2 the release checkpoint scores 15.48 exact pooled under the
a1-dark convention, ranging from 0 to 34.25 across the twenty factor cells, and 0.0 in every cell
under a1-light (Figure~\ref{fig:heatmap}, left).
% receipt: evidence/tools/reports/paper/real_photo_2026-09-11/a1_light_2026-10-04/score_cq_cpu_20261004.json (a1_dark_armA, a1_light_primary)
One real-board check points the same way. The same ten positions were photographed on one physical
board with a dark a1, and on a later day, with a different phone and photographer, with a light a1.
The release checkpoint scores 76.25 and 57.97 per-square on the two sets of ten, with no exact
board in either, and the light arm scores lower on 8 of 10 positions. The arms differ in more than
parity, so this is a descriptive check, not a paired probe.

The corrected statement is about per-square accuracy. The convention accounts for about two thirds
of the wrong squares, and most boards still fail exactly after the correction (18/192 and 50/288).
On this cohort the residual is family-associated: two families stay at 0/48 exact while another
reaches 19/48. Positions differ between families here, so this describes the unchanged release
checkpoint; it is not a paired geometry effect, and it says nothing about what adaptation later
fixes. The paired geometry cells of v2 and the adaptation runs are the place to read those
(Section~\ref{sec:adaptation}). The trap had been written in our own research notes before the
benchmark was built, and we missed it anyway. Making the convention an explicit factor let it be
measured instead of inferred, which is why v2 carries it as its fifth axis.

A bounded follow-up on the same 192 positions found that, under the standard convention, a
ground-truth square crop raised exact boards from 18 to 27 while lowering per-square accuracy from
85.24 to 84.24. Framing is therefore not irrelevant once the convention is corrected. The
interaction is a trade-off measured on development positions with privileged localisation, not a
deployable improvement.
